\SourceLeaf{523}\par\medskip\Needspace{4\baselineskip}\pdfbookmark[1]{372.}{sectio-octava-art-372}\centerline{\textbf{372.}}\nobreak\smallskip
If $n$ is an (odd) prime, the residues $m$ consist of zero, for which $M=1$, and $\frac12(n-1)$ other numbers, for which $M=2$. Denoting these residues (excluding the residue $0$), without further specification, by $\mu$, our progression is $=1+2\Sigma r^\mu$. Further, if $\nu$ denotes, without further specification, all remaining numbers below $n$, whose number is likewise $\frac12(n-1)$ and which comprise all quadratic nonresidues of $n$ below $n$, then evidently
\[
1+\Sigma r^\mu+\Sigma r^\nu=1+r+rr+r^3\dotsfour+r^{n-1}=\frac{1-r^n}{1-r}=0
\]
Thus, setting the sum of our progression, that is $1+2\Sigma r^\mu=A$, we have $1+2\Sigma r^\nu=-A$, and also $\Sigma r^\mu-\Sigma r^\nu=A$.
\par By Article $356$ we therefore have $A=\pm\surd n$ or $\pm\surd{-n}$, according as $n$ is $\equiv1$ or $\equiv3\, (\mathrm{mod}.4)$. But this does not yet determine the sign of the root.
\par If, in our progression, which we shall denote by $\Pi$, we replace $r$ by another similar root of the equation $x^n-1=0$, say $r'=r^\lambda$, let $\Pi'$ be the result.
\par\medskip\Needspace{26\baselineskip}\pdfbookmark[1]{373.}{sectio-octava-art-373}\centerline{\textbf{373.}}\nobreak\smallskip
If $n$ is the square or a higher power of a prime, say $=p^\pi$, some residues $m$ consist of numbers not divisible by $p$; others are divisible by $pp$ but not by a higher power of $p$; others can be divided by $p^4$ but not. by $p^5$; and so on to those divisible by $p^{\pi-2}$ but not by $p^{\pi-1}$, or by $p^{\pi-1}$ but not by $p^\pi$, according as $\pi$ is even or odd. Finally there is the residue $0$, the only one divisible by $p^\pi$ (compare Article $102$). Now denote, without further specification, by $\mu$ the quadratic residues of the number $p$ below $p$, excluding zero (whose number is $=\frac12(p-1)$). Those different classes of residues are then represented as follows. The first, which are not divisible by $p$, are represented by $\mu+kp$, where for $k$ we must substitute all integers from $0$ to $p^{\pi-1}-1$, so that the number of all residues contained in this form is $=\frac12(p-1)p^{\pi-1}$; for each of these, $M=2$. The sum of all terms in $\Pi$ corresponding to these residues is
\[
\begin{aligned}
&=2\Sigma r^{\mu+kp}=2\Sigma r^\mu\pdot\Sigma r^{kp}\\
&=2\Sigma r^\mu\pdot\frac{r^{p^\pi}-1}{r^p-1}=0
\end{aligned}
\]
The second class of residues is represented by $\mu pp+kp^3$, where for $k$ we must substitute all integers from $0$ to $p^{\pi-3}-1$, so that the number of all residues contained in this
